\subsection{Union and Intersection of Sets} 
We can combine sets in various ways.
\begin{definition}[Union]
The \term{union} of two sets is the set of all objects that are in at least one of the sets. We write:

\noindent\centering$A\cup B=\Set{x\mid x\in A\mbox{ or }x\in B}$
\end{definition}
\begin{Note}To be in the union, an object must be in one set or the other \emph{or both}. In math, we always mean an ``inclusive or.''\end{Note}
\begin{definition}[Intersection]
The \term{intersection} of two sets is the set of all objects that are in both sets. We write:

\noindent\centering$A\cap B=\Set{x\mid x\in A\mbox{ and }x\in B}$
\end{definition}
\begin{Example} Let $A=\Set{1,2,3,4,5}$ and $B=\Set{-2,2,4,6}$.\\[1ex]
$A\cup B=\Set{\makeblank}$\hfill
$A\cap B=\Set{\makeblank}$
\end{Example}
\begin{definition}[Empty Set]
A set which contains no objects is the \term{empty set}, which we represent with the symbol $\emptyset$.
\end{definition}
\begin{Example}Find each intersection below.\\[1ex]
\makeblankfill[0.6\textwidth]{$(-4,3]\cap\Z=$}\\[1ex]
\makeblankfill[0.6\textwidth]{$(-\infty,2)\cap[1,\infty)=$}\\[1ex]
\makeblankfill[0.6\textwidth]{$(0,1)\cap\Z=$}\\[1ex]
\makeblankfill[0.6\textwidth]{$(-3,-1)\cap(1,3)=$}
\end{Example}
\begin{Example}Find each union below.\\[1ex]
\makeblankfill[0.6\textwidth]{$[-3,4]\cup (2,6)=$}\\[1ex]
\makeblankfill[0.6\textwidth]{$[-8,\infty)\cup (2,5)=$}\\[1ex]
\makeblankfill[0.6\textwidth]{$(-\infty,2)\cup [1,\infty)=$}\\[1ex]
\makeblankfill[0.6\textwidth]{$(-\infty,-4]\cup (1,3)=$}
\end{Example}
\begin{itemize}
\item The intersection of two intervals is either an interval or empty.
\item The union of two intervals may be an interval, all real numbers, or neither.
\item If the union of two intervals contains a ``gap,'' we leave it as-is.
\end{itemize}